\section{Introduction}
\label{sec:introduction}

The intersection of physics-based modeling and machine learning has emerged as a frontier for scientific prediction, particularly in domains where first-principles understanding exists but computational or modeling limitations prevent accurate forecasting. Climate science presents a compelling testbed for such hybrid approaches: while many governing physical processes are understood to the first order, comprehensive dynamical models remain imperfect, and limited observational records mean that prediction often relies on climate-model simulations rather than purely data-driven machine learning. El Ni\~no-Southern Oscillation (ENSO), the dominant mode of inter-annual climate variability~\cite{mcphaden2006enso,timmermann2018enso}, exemplifies these challenges. Despite decades of research and operational forecasting efforts, predicting ENSO beyond six months remains difficult due to chaotic dynamics, model biases, and the so-called ``spring predictability barrier''~\cite{webster1992monsoon,duan2013predictability} where forecast skill drops sharply when predictions must traverse the boreal spring season.

Traditional approaches to ENSO forecasting face many challenges. First, dynamical models based on coupled ocean-atmosphere general circulation models (GCMs)~\cite{flato2013evaluation} encode comprehensive physical processes but suffer from systematic biases, require enormous computational resources, and struggle to outperform simpler statistical methods at lead times beyond one season~\cite{barnston2012skill,Latif1998JGR, Jin2008ClimDyn}. Second, deep learning and statistical models exploit observed relationships between climate variables but typically lack interpretability and may fail when climate statistics shift---a concerning limitation given anthropogenic climate change \cite{Dmitri_model, TangangTangMonahanHsieh1998NNEOF,UbilavaHelmers2013STAR, PenlandMagorian1993LIM}.
Recent advances in deep learning have prompted efforts to apply neural networks directly to climate prediction, with notable successes in multi-year ENSO forecasting~\cite{ham2019deep, Hu2021ResCNNENSO, Yan2020TCNENSO, ZhouZhang2023GeoformerENSO, Mu2024ENSOPhyNet}, albeit so far only in hindcasts and not in real-time. However, these purely data-driven approaches face a fundamental obstacle: the observational record provides at most a few hundred independent samples for training, far fewer than the millions typically required for deep learning to excel. These models have to rely on climate model simulation data with careful choices of the datasets and modalities to use, usually resulting in formidable storage and computational resources. On top of the resource challenges, these models often lack explainability, which is highly valuable for understanding natural phenomena.

This data scarcity and lack of explainability motivate the development of hybrid physics-machine learning models that combine the interpretability and sample efficiency of physics-based approaches with the flexibility and representational power of neural networks~\cite{willard2023integrating,yin2021augmenting}. The central question we address is: \emph{How can we best integrate physical structure with learned components for climate forecasting when data is severely limited?}

We build upon the Extended Recharge Oscillator (XRO) model~\cite{zhao2024xro}, a recent physics-based framework that achieves state-of-the-art ENSO forecast skill (better than the real-time skill of International Research Institute for Climate and Society (IRI) operational models 
and comparable to the most skilful AI forecast model ~\cite{ZhouZhang2023GeoformerENSO,Ham2019DeepLearningENSO}), 
by encoding the fundamental recharge-discharge dynamics~\cite{jin1997equatorial} of the tropical Pacific coupled with seasonal modulation and polynomial nonlinearities. Our contribution is the Neural Residual Extended Recharge Oscillator (NXRO) framework, inspired by neural ODE ~\cite{chen2018neuralode,rubanova2019latent}, physics-informed machine learning~\cite{karniadakis2021physics}, and differentiable physics~\cite{um2020solver,list2022learned}. This decomposition explicitly separates dynamics into a linear physics-based component and a nonlinear correction, providing strong inductive biases that enable learning from limited data. By elucidating the mathematical formulation of the base XRO, we develop 43 model variants spanning from pure physics (XRO) to pure machine learning, enabling systematic investigation of how much physical structure is necessary and where neural flexibility provides benefit.

We demonstrate the effectiveness of our framework via both out-of-sample deterministic and stochastic forecasts. Our results on deterministic forecasts highlight the importance of both physical inductive bias and the necessity of end-to-end optimization. Interestingly, we also find that adding more data samples (through climate model simulation) does not necessarily translate to better model performance, pointing to the fundamental bias in climate model simulation against observational record for time series forecasting tasks. For stochastic forecast, we introduce a simple two-stage stochastic training procedure, which achieves 13\% improvement in Continuous Ranked Probability Score (CRPS)~\cite{hersbach2000,gneiting2007} over XRO and produces better-calibrated uncertainty estimates. 